\documentclass[UTF8,11pt]{ctexart}
\usepackage[a4paper,margin=22mm]{geometry}
\usepackage{graphicx,enumitem,amssymb,tabularx}
\usepackage[normalem]{ulem}
\usepackage{xeCJKfntef}
\setlength{\parindent}{0pt}
\setlength{\parskip}{.6em}
\sloppy
\begin{document}
\section*{2015年12月英语四级真题(第3套)}
\textbf{2015} 年\textbf{12} 月大学英语四级考试真题\textbf{(}第\textbf{3} 套\textbf{)}


\subsection*{\textbf{Part} I\textbf{Writing ( 30 minutes)}}


Directions: For this part, you are allowed 30 minutes to write an essay commenting on the saying" Never go out there to see what happens, go out there to make things happen." You can cit examples to illustrate the importance of being participants rather than mere onlookers in life. You should write at least 120 words but no more than 180 words.


\subsection*{\textbf{Part} II\textbf{Listening Comprehension ( 30 minutes)}}


Sheet 1 with a single line through the centre. 1 A. Children should be taught to be more careful.


\begin{itemize}[leftmargin=2em,itemsep=.2em,topsep=.4em]\item[B.] Children shouldn't drink so much orange juice.
\item[C.] There is no need for the man to make such a fuss.
\item[D.] Timmy should learn to do things in the right way.
\end{itemize}

\noindent\begin{tabularx}{\linewidth}{@{}lXXXX@{}}2. & \textbf{A.} Fitness training. & \textbf{B.} The new job offer. & \textbf{C.} Computer programming. & \textbf{D.} Directorship of the club.\\\end{tabularx}\par

3.


\begin{itemize}[leftmargin=2em,itemsep=.2em,topsep=.4em]\item[A.] He needs to buy a new sweater.
\item[B.] He has got to save on fuel bills.
\item[C.] The fuel price has skyrocketed.
\item[D.] The heating system doesn't work.
\end{itemize}

\noindent\begin{tabularx}{\linewidth}{@{}lXXXX@{}}4. & \textbf{A.} Committing theft. & \textbf{B.} Taking pictures. & \textbf{C.} Window shopping. & \textbf{D.} Posing for the camera.\\\end{tabularx}\par

5.


\begin{itemize}[leftmargin=2em,itemsep=.2em,topsep=.4em]\item[A.] She is taking some medicine.
\item[B.] She has not seen a doctor yet.
\item[C.] She does not trust the man's advice.
\item[D.] She has almost recovered from the cough.
\end{itemize}

6.


\begin{itemize}[leftmargin=2em,itemsep=.2em,topsep=.4em]\item[A.] Pamela's report is not finished as scheduled.
\item[B.] Pamela has a habit of doing things in a hurry.
\item[C.] Pamela is not good at writing research papers.
\item[D.] Pamela's mistakes could have been avoided.
\end{itemize}

\noindent\begin{tabularx}{\linewidth}{@{}lXXXX@{}}7. & \textbf{A.} In the left-luggage office. & \textbf{B.} At the hotel reception. & \textbf{C.} In a hotel room. & \textbf{D.} At an airport.\\\end{tabularx}\par

8.


\begin{itemize}[leftmargin=2em,itemsep=.2em,topsep=.4em]\item[A.] She was an excellent student at college.
\item[B.] She works in the entertainment business.
\item[C.] She is fond of telling stories in her speech.
\item[D.] She is good at conveying her message.
\end{itemize}

\textbf{Questions 9 to 11 are based on the conversation you have just heard.}


9.


\begin{itemize}[leftmargin=2em,itemsep=.2em,topsep=.4em]\item[A.] Arranging the woman's appointment with Mr. Romero.
\item[B.] Fixing the time for the designer's latest fashion show.
\item[C.] Talking about an important gathering on Tuesday.
\item[D.] Preparing for the filming on Monday morning.
\end{itemize}

10.


\begin{itemize}[leftmargin=2em,itemsep=.2em,topsep=.4em]\item[A.] Her travel to Japan.
\item[B.] The awards ceremony.
\item[C.] The proper hairstyle for her new role.
\item[D.] When to start the make-up session.
\end{itemize}

11.


\begin{itemize}[leftmargin=2em,itemsep=.2em,topsep=.4em]\item[A.] He is Mr. Romero's agent.
\item[B.] He is an entertainment journalist.
\item[C.] He is the woman's assistant.
\item[D.] He is a famous movie star.
\end{itemize}

\textbf{Questions 12 to 15 are based on the conversation you have just heard.}


12.


\begin{itemize}[leftmargin=2em,itemsep=.2em,topsep=.4em]\item[A.] Make an appointment for an interview.
\item[B.] Send in an application letter.
\item[C.] Fill in an application form.
\item[D.] Make a brief self-introduction on the phone.
\end{itemize}

13.


\begin{itemize}[leftmargin=2em,itemsep=.2em,topsep=.4em]\item[A.] Someone having a college degree in advertising.
\item[B.] Someone experienced in business management.
\item[C.] Someone ready to take on more responsibilities.
\item[D.] Someone willing to work beyond regular hours.
\end{itemize}

\noindent\begin{tabularx}{\linewidth}{@{}lXXXX@{}}14. & \textbf{A.} Travel opportunities. & \textbf{B.} Handsome pay. & \textbf{C.} Prospects for promotion. & \textbf{D.} Flexible working hours.\\\end{tabularx}\par

15.


\begin{itemize}[leftmargin=2em,itemsep=.2em,topsep=.4em]\item[A.] It depends on the working hours.
\item[B.] It is about 500 pounds a week.
\item[C.] It will be set by the Human Resources.
\item[D.] It is to be negotiated.
\end{itemize}

\subsection*{\textbf{Section B}}


\textbf{Directions..}In this section, you will hear 3 short passages. At the end of each passage, you will hear some questions. Both the passage and the questions will be spoken only once. After you hear a question, you must choose the best answer from the four choices marked A., B, C and D..Then mark the corresponding letter on Answer Sheet 1 with a single line through the center.


\subsection*{Passage One}


\textbf{Questions 16 to 19 are based on the passage you have just heard.}


16.


\begin{itemize}[leftmargin=2em,itemsep=.2em,topsep=.4em]\item[A.] To give customers a wider range of choices.
\item[B.] To make shoppers see as many items as possible.
\item[C.] To supply as many varieties of goods as it can.
\item[D.] To save space for more profitable products.
\end{itemize}

17.


\begin{itemize}[leftmargin=2em,itemsep=.2em,topsep=.4em]\item[A.] On the top shelves.
\item[B.] On the bottom shelves.
\item[C.] On easily accessible shelves.
\item[D.] On clearly marked shelves.
\end{itemize}

18.


\begin{itemize}[leftmargin=2em,itemsep=.2em,topsep=.4em]\item[A.] Many of them buy things on impulse.
\item[B.] A few of them are fathers with babies.
\item[C.] A majority of them are young couples.
\item[D.] Over 60\% of them make shopping lists.
\end{itemize}

19.


\begin{itemize}[leftmargin=2em,itemsep=.2em,topsep=.4em]\item[A.] Sales assistants promoting high margin goods.
\item[B.] Sales assistants following customers around.
\item[C.] Customers competing for good bargains.
\item[D.] Customers losing all sense of time.
\end{itemize}

\subsection*{\textbf{Passage Two}}


\textbf{Questions 20 to 22 are based on the passage you have just heard.}


20.


\begin{itemize}[leftmargin=2em,itemsep=.2em,topsep=.4em]\item[A.] Teaching mathematics at a school.
\item[B.] Doing research in an institute.
\item[C.] Studying for a college degree.
\item[D.] Working in a hi-tech company
\end{itemize}

21.


\begin{itemize}[leftmargin=2em,itemsep=.2em,topsep=.4em]\item[A.] He studied the designs of various clocks.
\item[B.] He did experiments on different materials.
\item[C.] He bought an alarm clock with a pig face.
\item[D.] He asked different people for their opinions.
\end{itemize}

22.


\begin{itemize}[leftmargin=2em,itemsep=.2em,topsep=.4em]\item[A.] Its automatic mechanism.
\item[B.] Its manufacturing process.
\item[C.] Its way of waking people up.
\end{itemize}

\begin{itemize}[leftmargin=2em,itemsep=.2em,topsep=.4em]\item[D.] Its funny-looking pig face.
\end{itemize}

\subsection*{\textbf{Passage Three}}


\textbf{Questions 23 to 25 are based on the passage you have just heard.}


23.


\begin{itemize}[leftmargin=2em,itemsep=.2em,topsep=.4em]\item[A.] It is often caused by a change of circumstances.
\item[B.] It actually doesn't require any special treatment.
\item[C.] It usually appears all of a sudden.
\item[D.] It generally lasts for several years.
\end{itemize}

24.


\begin{itemize}[leftmargin=2em,itemsep=.2em,topsep=.4em]\item[A.] They cannot mix well with others.
\item[B.] They irrationally annoy their friends.
\item[C.] They depend heavily on family members.
\item[D.] They blame others for ignoring their needs.
\end{itemize}

25.


\begin{itemize}[leftmargin=2em,itemsep=.2em,topsep=.4em]\item[A.] They lack consistent support from peers.
\item[B.] They doubt their own popularity.
\item[C.] They were born psychologically weak.
\item[D.] They focus too much on themselves.
\end{itemize}

\subsection*{\textbf{Part} III\textbf{Reading Comprehension(40 minutes)}}


\subsection*{Section A}


Directions: In this section, there is a passage with ten blanks. You are required to select one word for each blank from a list of choices given, in a word bank following the passage. Read the passage through carefully before making your choices. Each. choice in the bank is identified by a letter. Please mark the corresponding letter for each item on Answer Sheet2 with a single line through the center. You may not use any of the words in the bank more than once.


\textbf{Questions 36 to 45 are based on the following passage.}


Children do not think the way adults do. For most of the first year of life, if something is out of sight, it's out of mind. If you cover a baby's\uline{36}toy with a piece of cloth, the baby thinks the toyhas disappeared and stops looking for it. A 4-year-old may\uline{ 37 }that a sister has more fruit juice when it is only the shapes of the glasses that differ, not the\uline{38 }of juice.


Yet children are smart in their own way. Like good little scientists, children are always testing their child-sized\uline{39 }about how things work. When your child throws her spoon on the floor for the sixth time as you try to feed her, and you say, "That's enough! I will not pick up your spoon again!" the child will\uline{ 40 }test your claim. Are you serious? Are you angry? What will happen if she throws the spoon again? She is not doing this to drive you\uline{41}; rather, she is learning that her desires and yours can differ, and that sometimes those\uline{42 }are important and sometimes they are not.


How and why does children's thinking change? In the 1920s, Swiss psychologist Jean Piaget proposed that children's cognitive (认知的) abilities unfold\uline{ 43}, like the blooming of a flower, almost independent of what else is\uline{44}in their lives. Although many of his specific conclusions havebeen\uline{45 }or modified over the years, his ideas inspired thousands of studies by investigators all over the world.


\begin{itemize}[leftmargin=2em,itemsep=.2em,topsep=.4em]\item[A.] advocate
\item[B.] amount
\item[C.] confirmed
\item[D.] crazy
\item[E.] definite
\item[F.] differences
\item[G.] favorite
\item[H.] happening
\item[I.] Immediately
\item[J.] Naturally
\item[K.] Obtaining
\item[L.] Primarily
\item[M.] Protest
\item[N.] Rejected
\item[O.] theories
\end{itemize}

\subsection*{\textbf{Section B}}


\textbf{Directions:} In this section, you are going to read a passage with, ten statements attached to it. Each statement contains information given in one of the paragraphs. Identify the paragraph from which the information is derived. You may choose a paragraph more than once. Each paragraph is marked with a letter. Answer the questions by marking the corresponding letter on Answer Sheet 2.The Perfect Essay


\begin{itemize}[leftmargin=2em,itemsep=.2em,topsep=.4em]\item[A.] Looking back on too many years of education, I can identify one truly impossible teacher. She
\end{itemize}

cared about me, and my intellectual life, even when I didn't. Her expectations were high--impossibly so. She was an English teacher. She was also my mother.


\begin{itemize}[leftmargin=2em,itemsep=.2em,topsep=.4em]\item[B.] When good students turn in an essay, they dream of their instructor returning it to them in
\end{itemize}

exactly the same condition, save for a single word added in the margin of the final page : "Flawless." This dream came true for me one afternoon in the ninth grade. Of course, I had heard that genius could show itself at an early age, so I was only slightly taken aback that I had achieved perfection at the tender age of 14.Obviously, I did what any professional writer would do; I hurried off to spread the good news. I didn't get very far. The first person I told was my mother.


\begin{itemize}[leftmargin=2em,itemsep=.2em,topsep=.4em]\item[C.] My mother, who is just shy of five feet tall, is normally incredibly soft-spoken, but on the rare
\end{itemize}

occasion when she got angry, she was terrifying. I am not sure if she was more upset by my hubris(得意忘形) or by the fact that my English teacher had let my ego get so out of hand. In any event, my mother and her red pen showed me how deeply flawed a flawless essay could be. At the time, I am sure she thought she was teaching me about mechanics, transitions (过渡), structure, style and voice. But what I learned, and what stuck with me through my time teaching writing at Harvard, was a deeper lesson about the nature of creative criticism.


\begin{itemize}[leftmargin=2em,itemsep=.2em,topsep=.4em]\item[D.] First off, it hurts. Genuine criticism, the type that leaves a lasting mark on you as a writer, also
\end{itemize}

leaves an existential imprint (印记) on you as a person. I have heard people say that a writer should never take criticism personally. I say that we should never listen to these people.


\begin{itemize}[leftmargin=2em,itemsep=.2em,topsep=.4em]\item[E.] Criticism, at its best, is deeply personal, and gets to the heart of why we write the way we do.
\end{itemize}

The intimate nature of genuine criticism implies something about who is able to give it, namely, someone who knows you well enough to show you how your mental life is getting in the way of good writing. Conveniently, they are also the people who care enough to see you through this painful realization. For me it took the form of my first, and I hope only, encounter with writer's block--I was not able to produce anything for three years.


\begin{itemize}[leftmargin=2em,itemsep=.2em,topsep=.4em]\item[F.] Franz Kafka once said: "Writing is utter solitude (独处), the descent into the cold abyss (深渊)
\end{itemize}

of oneself." My mother's criticism had shown me that Kafka is right about the cold abyss, and when you make the introspective (内省的) descent that writing requires you are not always pleased by what you find. But, in the years that followed, her sustained tutoring suggested that Kafka might be wrong about the solitude. I was lucky enough to find a critic and teacher who was willing to make the journey of writing with me."It is a thing of no great difficulty," according to Plutarch, "to raise objections against another man's speech, it is a very easy matter; but to produce a better in its place is a work extremely troublesome." I am sure I wrote essays in the later years of high school without my mother's guidance, but I can't recall them. What I remember, however, is how she took up the "extremely troublesome" work of ongoing criticism.


\begin{itemize}[leftmargin=2em,itemsep=.2em,topsep=.4em]\item[G.] There are two ways to interpret Plutarch when he suggests that a critic should be able to
\end{itemize}

produce "a better in its place." In a straightforward sense, he could mean that a critic must be more talented than the artist she critiques (评论).My mother was well covered on this count. But perhaps


Plutarch is suggesting something slightly different, something a bit closer to Marcus Cicero's claim that one should "criticize by creation, not by finding fault." Genuine criticism creates a precious opening for an author to become better on his own terms--a process that is often extremely painful, but also almost always meaningful.


\begin{itemize}[leftmargin=2em,itemsep=.2em,topsep=.4em]\item[H.] My mother said she would help me with my writing, but first I had to help myself. For each
\end{itemize}

assignment, I was to write the best essay I could. Real criticism is not meant to find obvious mistakes, so if she found any--the type I could have found on my own--I had to start from scratch. From scratch. Once the essay was "flawless," she would take an evening to walk me through my errors. That was when true criticism, the type that changed me as a person, began.


\begin{itemize}[leftmargin=2em,itemsep=.2em,topsep=.4em]\item[I.] She criticized me when I included little-known references and professional jargon (行话).She
\end{itemize}

hadno patience for brilliant but irrelevant figures of speech."Writers can't bluff (虚张声势) their way through ignorance." That was news to me--I would need to freed another way to structure my daily existence.


\begin{itemize}[leftmargin=2em,itemsep=.2em,topsep=.4em]\item[J.] She trimmed back my flowery language, drew lines through my exclamation marks and argued
\end{itemize}

for the value of restraint in expression."John," she almost whispered. I leaned in to hear her: "I can't hear you when you shout at me." So I stopped shouting and bluffing, and slowly my writing improved.


\begin{itemize}[leftmargin=2em,itemsep=.2em,topsep=.4em]\item[K.] Somewhere along the way I set aside my hopes of writing that flawless essay. But perhaps I
\end{itemize}

missed something important in my mother's lessons about creativity and perfection. Perhaps the point of writing the flawless essay was not to give up, but to never willingly finish. Whit man repeatedly reworked "Song of Myself' between 1855 and 1891.Repeatedly.We do our absolute best with a piece of writing, and come as close as we can to the ideal. And, for the time being, we settle. In critique, however, we are forced to depart, to give up the perfection we thought we had achieved for the chance of being even a little bit better. This is the lesson I took from my mother: If perfection were possible, it would not be motivating.


46.The author was advised against the improper use of figures of speech.


47.The author's mother taught him a valuable lesson by pointing out lots of flaws in his seemingly


perfect essay.


48.A writer should polish his writing repeatedly so as to get closer to perfection.


49.Writers may experience periods of time in their life when they just can't produce anything.


50.The author was not much surprised when his school teacher marked his essay as "flawless".


51.Criticizing someone's speech is said to be easier than coming up with a better one.


52.The author looks upon his mother as his most demanding and caring instructor.


53.The criticism the author received from his mother changed him as a person.


54.The author gradually improved his writing by avoiding fancy language.


55.Constructive criticism gives an author a good start to improve his writing.


\subsection*{\textbf{Section C}}


\textbf{Directions:} There are 2 passages in this section. Each passage is followed by some questions or unfinished statements. For each of them there are four choices marked A, B, C. and D .You should decide on the best choice and mark the corresponding letter on Answer Sheet 2 with a single line through the center.


\subsection*{\textbf{Passage One}}


\textbf{Questions 56 to 60 are based on the following passage.}


Could you reproduce Silicon Valley elsewhere, or is there something unique about it?


It wouldn't be surprising if it were hard to reproduce in other countries, because you couldn't reproduce it in most of the US either. What does it take to make a Silicon Valley?


It's the right people. If you could get the right ten thousand people to move from Silicon Valley to Buffalo, Buffalo would become Silicon Valley.


You only need two kinds of people to create a technology hub (中心) : rich people and nerds (痴迷科研的人).


Observation bears this out. Within the US, towns have become star, up hubs if and only if they have both rich people and nerds. Few startups happen in Miami, for example, because although it's full of rich people, it has few nerds. It's not the kind of place nerds like.


Whereas Pittsburgh has the opposite problem: plenty of nerds, but no rich people.The top US Computer Science departments are said to be MIT, Stanford, Berkeley, and Carnegie-Mellon. MIT yielded Route 128.Stanford and Berkeley yielded Silicon Valley. But what did Carnegie-Mellon yield in Pittsburgh? And what happened in Ithaca, home of Cornell University, which is also high on the list?


I grew up in Pittsburgh and went to college at Cornell, so I can answer for both. The weather is terrible, particularly in winter, and there's no interesting old city to make up for it, as there is in Boston. Rich people don't want to live in Pittsburgh or Ithaca. So while there're plenty of hackers (电脑迷) who could start startups, there's no one to invest in them


Do you really need the rich people? Wouldn't it work to have the government invest in the nerds?


No, it would not. Start up investors are a distinct type of rich people. They tend to have a lot of experience themselves in the technology business. This helps them pick the right startups, and means they can supply advice and connections as well as money. And the fact that they have a personal stake in the outcome makes them really pay attention.


56.What do we learn about Silicon Valley from the passage?


\begin{itemize}[leftmargin=2em,itemsep=.2em,topsep=.4em]\item[A.] Its success is hard to copy anywhere else.
\item[B.] It is the biggest technology hub in the US.
\item[C.] Its fame in high technology is incomparable.
\item[D.] It leads the world in information technology.
\end{itemize}

57.What makes Miami unfit to produce a Silicon Valley?


\begin{itemize}[leftmargin=2em,itemsep=.2em,topsep=.4em]\item[A.] Lack of incentive for investment.
\item[B.] Lack of the right kind of talents.
\item[C.] Lack of government support.
\item[D.] Lack of famous universities.
\end{itemize}

58.In what way is Carnegie-Mellon different from Stanford, Berkeley and MIT?


\begin{itemize}[leftmargin=2em,itemsep=.2em,topsep=.4em]\item[A.] Its location is not as attractive to rich people.
\item[B.] Its science departments are not nearly as good.
\item[C.] It does not produce computer hackers and nerds.
\item[D.] It does not pay much attention to business startups.
\end{itemize}

59.What does the author imply about Boston?


\begin{itemize}[leftmargin=2em,itemsep=.2em,topsep=.4em]\item[A.] It has pleasant weather all year round.
\item[B.] It produces wealth as well as high-tech.
\item[C.] It is not likely to attract lots of investors and nerds.
\item[D.] It is an old city with many sites of historical interest.
\end{itemize}

60.What does the author say about startup investors?


\begin{itemize}[leftmargin=2em,itemsep=.2em,topsep=.4em]\item[A.] They are especially wise in making investments.
\item[B.] They have good connections in the government.
\item[C.] They can do more than providing money.
\item[D.] They are rich enough to invest in nerds.
\end{itemize}

\subsection*{\textbf{Passage Two}}


\textbf{Questions 61 to 65 are based on the following passage.}


It's nice to have people of like mind around. Agreeable people boost your confidence and allow you to relax and feel comfortable. Unfortunately, that comfort can hinder the very learning that can expand your company and your career.


It's nice to have people agree, but you need conflicting perspectives to dig out the truth .If everyone around you has similar views, your work will suffer from confirmation bias (偏颇).


Take a look at your own network. Do your contacts share your point of view on most subjects? If yes, it's time to shake things up .As a leader, it can be challenging to create an environment in which people will freely disagree and argue, but as the saying goes: From confrontation comes brilliance.


It's not easy for most people to actively seek conflict. Many spend their lives trying to avoid arguments. There's no need to go out and find people you hate, but you need to do some self-assessment to determine where you have become stale in your thinking. You may need to start by encouraging your current network to help you identify your blind spots.


Passionate, energetic debate does not require anger and hard feelings to be effective. But it does require moral strength. Once you have worthy opponents, set some ground rules so everyone understands responsibilities and boundaries. The objective of this debating game is not to win but to get to the truth that will allow you to move faster, farther, and better.


Fierce debating can hurt feelings, particularly when strong personalities are involved. Make sure you check in with your opponents so that they are not carrying the emotion of the battles beyond the battlefield. Break the tension with smiles and humor to reinforce the idea that this is friendly discourse and that all are working toward a common goal.


Reward all those involved in the debate sufficiently when the goals are reached. Let your sparring partners (拳击陪练) know how much you appreciate their contribution. The more they feel appreciated, the more they'll be willing to get into the ring next time.


61.What happens when you have like-minded people around you all the while?


\begin{itemize}[leftmargin=2em,itemsep=.2em,topsep=.4em]\item[A.] It will help your company expand more rapidly.
\item[B.] It will create a harmonious working atmosphere.
\item[C.] It may prevent your business and career from advancing.
\item[D.] It may make you feel uncertain about your own decisions.
\end{itemize}

62.What does the author suggest leaders do?


\begin{itemize}[leftmargin=2em,itemsep=.2em,topsep=.4em]\item[A.] Avoid arguments with business partners.
\item[B.] Encourage people to disagree and argue.
\item[C.] Build a wide and strong business network.
\item[D.] Seek advice from their worthy competitors.
\end{itemize}

63.What is the purpose of holding a debate?


\begin{itemize}[leftmargin=2em,itemsep=.2em,topsep=.4em]\item[A.] To find out the truth about an issue.
\item[B.] To build up people's moral strength.
\item[C.] To remove misunderstandings.
\item[D.] To look for worthy opponents.
\end{itemize}

64.What advice does the author give to people engaged in a fierce debate?


\begin{itemize}[leftmargin=2em,itemsep=.2em,topsep=.4em]\item[A.] They listen carefully to their opponents' views.
\item[B.] They show due respect for each other's beliefs.
\item[C.] They present their views clearly and explicitly.
\item[D.] They take care not to hurt each other's feelings.
\end{itemize}

65.How should we treat our rivals after a successful debate?


\begin{itemize}[leftmargin=2em,itemsep=.2em,topsep=.4em]\item[A.] Try to make peace with them.
\item[B.] Try to make up the differences.
\item[C.] Invite them to the ring next time.
\item[D.] Acknowledge their contribution.
\end{itemize}

\textbf{Part} IV\textbf{Translation( 30 minutes )} Directions: For this part, you are allowed 30 minutes to translate a passage from Chinese into English. You should write your answer on Answer Sheet 2.


今年在长沙举行了一年一度的外国人汉语演讲比赛.这项比赛证明是促进中国和世界其他地区文化交流的好方法.它为世界各地的年轻人提供了更好地了解中国的机会.来自87 个国家共计 126 位选手聚集在湖南省省会参加了从7 月6 日到8 月5 日进行的半决赛和决赛.比赛并不是唯一的活动.选手们还有机会参观了中国其他地区的著名景点和历史名胜.

\end{document}